%Experiment 4
%Develop a LaTeX script to create a simple document that consists of 2 sections 
%[Section1, Section2], and a paragraph with dummy text in each section. And also 
%include header [title of document] and footer [institute name, page number] in 
%the document. 
\documentclass[12pt]{article} % Document type and font size
\usepackage{graphicx} % For image support
\usepackage{fancyhdr} % For custom header and footer
\usepackage{lipsum}  % For dummy text generation
\pagestyle{fancy}    % Activate fancy page style
\fancyhf{}            % Clear default header and footer
\fancyhead[C]{THE HINDUSTAN TIMES}
% Header: document title
\fancyfoot[L]{GNDEC,LUDHIANA}
% Footer left: institute name
\fancyfoot[C]{\thepage}
% Footer center: page number
\fancyfoot[R]{URN:2435157}
% Footer right: registration number
\begin{document}
% Begin document content
\section{Lorem Ipsum Dolor}
% Section 1
\lipsum[1][1-3]
% Dummy paragraph
\subsection{Sit Amet Consectetur}
% Subsection under Section 1
\lipsum[2][1-3]
% Dummy paragraph
\section{Adipiscing Elit Sed}
% Section 2
\lipsum[3][1-3]
% Dummy paragraph
\subsection{Eiusmod Tempor Incididunt}
% Subsection under Section 2
\lipsum[4][1-3]
% Dummy paragraph
\subsection{Ut Labore Et}
% Another subsection under Section 2
\lipsum[5][1-3]
% Dummy paragraph
\end{document}
% End of document




%Experiment 5
%Develop a LaTeX script to create a document that displays the sample 
%Abstract/Summary.  
\documentclass{article}                          % Use the standard article class
\usepackage{lipsum}                              % Package to generate dummy text
\title{Experiment-05: Sample Abstract}           % Document title
\author{Saksham}                                 % Author name
\date{}                                          % Empty date
\begin{document}	                             % Begin document content
\maketitle                                       % Generates the title, author, and date
\begin{abstract}
This is a sample abstract.
It provides a summary of the document, including the purpose, methodology, and
conclusions.
The abstract should be concise, usually not more than 200 words, and allow readers
to quickly understand the scope of the work.
\end{abstract}                                   % Abstract section
\section{Introduction}
\lipsum[1]                                       % Dummy text for the Introduction section
\section{Conclusion}
\lipsum[1]                                       % Dummy text for the Conclusion section
\end{document}                                   % End of the document






%Experiment 6
%Develop a LaTeX script to create the Certificate Page of the Report [Use suitable 
%commands to leave the blank spaces for user entry]. 
\documentclass[12pt,a4paper]{article}
\usepackage[a4paper,margin=1in]{geometry}
\usepackage{setspace} \usepackage{fancyhdr}
\usepackage{ragged2e} \usepackage{times}
\setstretch{1.25} \pagestyle{fancy}
\fancyhf{}
\fancyfoot[C]{\thepage}
\fancyhead[L]{\textbf{Guru Nanak Dev Engineering College, Ludhiana}}
\fancyhead[R]{\textbf{Department of CSE}}
\begin{document}
\begin{titlepage}
\begin{center}
{\Large \textbf{Guru Nanak Dev Engineering College, Ludhiana}}\\[0.2cm]
{\large (An Autonomous College under UGC Act)}\\[0.3cm]
\rule{\textwidth}{0.6pt}\\[1cm]
{\LARGE \textbf{CERTIFICATE}}\\[1.2cm]
\noindent
\begin{minipage}{0.9\textwidth}
\justifying
This is to certify that the project/report entitled
\textbf{"\underline{\hspace{6cm}}"}, 
\textbf{\underline{\hspace{5.5cm}}},
submitted by 
\textbf{Branch: Computer Science and Engineering}, bearing
University Roll No. \textbf{\underline{\hspace{4cm}}}, in partial
fulfillment of the requirements for the award of the degree of
\textbf{\underline{\hspace{6cm}}} to \textbf{Guru Nanak Dev
Engineering College (Autonomous), Ludhiana}, is a bonafide record
of the work carried out by him/her under my supervision and
guidance during the academic year
\textbf{\underline{\hspace{4cm}}}.
\end{minipage}
\vspace{2cm}
\begin{minipage}{0.47\textwidth}
\centering
\textbf{Signature of the Guide}\\[1cm]
\begin{tabular}{@{}l@{}}
Name: \underline{\hspace{6cm}} \\[0.5cm]
Designation: \underline{\hspace{5cm}} \\[0.5cm]
Department: \underline{\hspace{5cm}} \\
\end{tabular}
\end{minipage}
\hfill
\begin{minipage}{0.47\textwidth}
\centering
\textbf{Signature of the Head of Department}\\[1cm]
\begin{tabular}{@{}l@{}}
Name: \underline{\hspace{6cm}} \\[0.5cm]
Designation: \underline{\hspace{5cm}} \\[0.5cm]
Department: \underline{\hspace{5cm}} \\
\end{tabular}
\end{minipage}
\vfill
\noindent
\begin{minipage}{0.45\textwidth}
\textbf{Date:} \underline{\hspace{5cm}}
\end{minipage}
\hfill
\begin{minipage}{0.45\textwidth}
\raggedleft \textbf{Place:} \underline{\hspace{5cm}}
\end{minipage}
\end{center}
\end{titlepage}
\end{document}






%Experiment 7
%Table
\documentclass[a4paper,12pt]{article}  
\usepackage{geometry}  
\usepackage{fancyhdr}  
\usepackage{caption}  
\usepackage{multirow}  
\usepackage{array}  
\usepackage{booktabs}  
\geometry{margin=1in}  
\pagestyle{fancy}  
\fancyhf{} % Clear default headers/footers  
\fancyhead[L]{Experiment 07}  
\fancyfoot[l]{GNDEC}  
\fancyfoot[R]{{URN: 2435157}}   
\renewcommand{\headrulewidth}{1pt}  
\renewcommand{\footrulewidth}{1pt}  
\title{\textbf{Student Record}}  
\author{SAKSHAM}  
\date{}  
\begin{document}  
\maketitle  
\thispagestyle{fancy}   
\noindent \large  
\vspace{1em}  
\begin{table}[h!]  
\centering  
\caption{Marks Obtained by Students in Three Subjects}  
\renewcommand{\arraystretch}{1.3}  
\begin{tabular}{|c|c|c|c|c|c|}  
\hline  
\multirow{2}{*}{\textbf{S.No}} &  
\multirow{2}{*}{\textbf{URN}} &  
\multirow{2}{*}{\textbf{Student Name}} &  
\multicolumn{3}{c|}{\textbf{Marks}} \\  
\cline{4-6}  
 & & & \textbf{Data Structure} & \textbf{OOPs} & \textbf{Maths III} \\  
\hline  
1 & 243518501 & Simran   & 89 & 60 & 90 \\  
\hline  
2 & 243518502 & Aman & 78 & 45 & 98 \\  
\hline  
3 & 243518503 & Divya  & 67 & 55 & 59 \\  
\hline  
\end{tabular}  
\end{table}  
\end{document} 







%Experiment 8
%side by side figures
\documentclass[10pt,a4paper]{article}  
\usepackage[utf8]{inputenc}  
\usepackage{graphicx}      
\usepackage{subcaption}     
\usepackage{fancyhdr}       
\usepackage[left=2.5cm,right=2.5cm,top=2.5cm,bottom=2.5cm]{geometry}  
\pagestyle{fancy}  
\fancyhf{}  
\fancyhead[L]{\textbf{Experiment 08}}  
\fancyhead[R]{By Saksham}  
\fancyfoot[L]{GNDEC}  
\fancyfoot[R]{URN: 2435157}  
\renewcommand{\headrulewidth}{1pt}  
\renewcommand{\footrulewidth}{1pt}  
\setlength{\headheight}{15pt}  
\setlength{\headsep}{25pt}  
\setlength{\footskip}{30pt}  
\begin{document}  
\title{{\LARGE \textbf{Experiment 08: Demonstration of Subfigures in \LaTeX}}}  
\author{SAKSHAM}  
\date{}  
\maketitle  
\thispagestyle{fancy}     
\vspace{6pt}  
\begin{center}  
{\large Understanding the use of subfigure environment for arranging images side
by-side} \end{center}  
\vspace{12pt}  
\noindent  
\textbf{Objective:}    
The aim of this experiment is to demonstrate how multiple related images can be 
arranged neatly within a single figure environment using the \texttt{subfigure} package 
in \LaTeX. This approach helps in presenting visual comparisons and grouped 
graphical outputs more effectively.  
\vspace{6pt}  
\textbf{Description:}    
In this experiment, three different mathematical graphs are displayed side-by-side. 
Each graph has its own caption and label, making it easy to compare their shapes and 
behaviors. Subfigures are especially useful in documentation, reports, and research 
papers where grouped visual information is required.  
\section*{Graphs Included}  
\begin{itemize}  
    \item \textbf{$y = x$} — A straight line representing a linear relationship.  
    \item \textbf{$y = 3\sin x$} — A sine wave with an amplitude of 3, showing periodic 
oscillations.  
    \item \textbf{$y = \frac{5}{x}$} — A rational function representing a hyperbola with 
inverse variation.  
\end{itemize}  
\section*{Subfigure Demo}  
\begin{figure}[h]  
    \centering  
    \begin{subfigure}[b]{0.3\textwidth}  
        \centering  
        \includegraphics[width=\textwidth]{graph1.png}  
        \caption{$y = x$}  
        \label{fig:y_equals_x}  
    \end{subfigure}  
    \hfill  
    \begin{subfigure}[b]{0.3\textwidth}  
        \centering  
        \includegraphics[width=\textwidth]{graph2.png}  
        \caption{$y = 3\sin x$}  
        \label{fig:three_sin_x}  
    \end{subfigure}  
    \hfill  
    \begin{subfigure}[b]{0.3\textwidth}  
        \centering  
        \includegraphics[width=\textwidth]{graph3.png}  
        \caption{$y = \frac{5}{x}$}  
        \label{fig:five_over_x}  
    \end{subfigure}  
    \caption{Three mathematical graphs arranged side-by-side using subfigures}  
    \label{fig:three_graphs}  
\end{figure}  
\end{document}







%Experiment 9
%Develop a LaTeX script to create a document that consists of two mathematical 
%equations.
\documentclass[12pt,a4paper]{article}  
\usepackage{geometry}  
\usepackage{amsmath, amssymb}  
\geometry{margin=1in}  
\usepackage{fancyhdr}          
\pagestyle{fancy}  
% for header and footer  
\fancyhead[R]{By Saksham} % clear all header fields  
\fancyhead[L]{\textbf{Experiment 09}}  
\fancyfoot{} % clear all footer fields  
\fancyfoot[L]{GNDEC} % Left footer  
\fancyfoot[R]{URN: 2435157} % Right footer  
\renewcommand{\footrulewidth}{1pt}   
\title{\textbf{Mathematical Equations in LaTeX}}  
\date{}  
\begin{document}  
\maketitle  
\thispagestyle{fancy} % Show header/footer on title page  
\section*{Mathematical Equations}  
This document demonstrates two mathematical equations with brief explanations 
written in \LaTeX.  
\subsection*{1. Euler’s Formula}  
Euler’s formula is a fundamental equation in complex analysis. It beautifully connects 
the exponential function with the trigonometric functions sine and cosine:  
\begin{equation} 
e^{i\theta} 
\cos\theta + i\sin\theta  
\end{equation}  
= 
This equation shows how complex exponentials can represent rotations in the complex 
plane. It also forms the basis of Euler's identity, which is often called the most elegant 
equation in mathematics.  
\vspace{0.3cm}  
\subsection*{2. Quadratic Formula}  
The quadratic formula provides the solutions to the quadratic equation \(ax^2 + bx + c 
= 0\). It is derived by completing the square and is used to find the roots of any second
degree polynomial:  
\begin{equation} x = \frac{-b \pm 
\sqrt{b^2 - 4ac}}{2a}  
\end{equation}  
This formula is widely used in algebra, physics, and engineering to find real or complex 
solutions. The discriminant (\(b^2 - 4ac\)) helps determine whether the roots are real, 
equal, or complex.  
  
\end{document}










%Experiment 10
%Develop a LaTeX script to demonstrate the presentation of Numbered theorems, 
%definitions, corollaries, and lemmas in the document.  
\documentclass{article}  
\usepackage{geometry}  
\geometry{margin=1in}  
\usepackage{fancyhdr}  
\pagestyle{fancy}  
\fancyhf{} 
\fancyhead[R]{By Saksham}
\fancyhead[L]{Experiment 10}  
\fancyfoot[L]{GNDEC}  
\fancyfoot[R]{URN: 2435157}  
\renewcommand{\headrulewidth}{1pt} 
\renewcommand{\footrulewidth}{1pt}  
\newtheorem{theorem}{Theorem}[section]  
\newtheorem{definition}[theorem]{Definition}  
\newtheorem{corollary}[theorem]{Corollary}  
\newtheorem{lemma}[theorem]{Lemma}  
\title{\LARGE \textbf{Theorem EnvironmentsX}}  
\author{SAKSHAM}  
\date{}  
\begin{document}  
\maketitle  
\thispagestyle{fancy}  
\section{Introduction}  
\large  
This document demonstrates how to present mathematical statements using LaTeX's 
theorem environments.  
\section{Main Results}  
\large  
\begin{theorem}[Pythagorean Theorem]  
In a right-angled triangle, the square of the hypotenuse equals the sum of the squares 
of the other two sides:  
\begin{equation} a^2 
+ b^2 = c^2  
\label{eq:pythagoras}  
\end{equation}  
\end{theorem}  
\begin{proof}  
\large  
Let the triangle have sides \(a\), \(b\), and hypotenuse \(c\). By constructing squares 
on each side and comparing areas, we find that the area of the square on the 
hypotenuse equals the sum of the areas on the other two sides.  
Equation~\ref{eq:pythagoras} confirms this relationship.  
\renewcommand{\qedsymbol}{}  

\end{proof}  
\begin{definition}\large  
A prime number is a natural number greater than 1 that has no positive divisors other 
than 1 and itself.  
\end{definition}  
\begin{definition}[Fibration]\large  
A fibration is a continuous map between topological spaces that satisfies the homotopy 
lifting property.  
\end{definition}  
\begin{corollary}\large  
The number 2 is the only even prime number.  
\end{corollary}  
\begin{lemma}\large  
If \(a\) divides \(b\) and \(b\) divides \(c\), then \(a\) divides \(c\).  
\end{lemma}  
\section{Conclusion}\large  
We have illustrated the use of theorem environments in LaTeX with proper numbering, 
structure, and equation labeling.  
\end{document} 

%Experiment 11
%Develop a LaTeX script to create a document that consists of two paragraphs 
%with a minimum of 10 citations in it and display the reference in the section. 
\documentclass{article}  
\usepackage{cite}  
\usepackage{fancyhdr}  
\pagestyle{fancy}  
\fancyhf{}
\fancyhead[R]{By Saksham}
\fancyhead[L]{Experiment 11}  
\fancyfoot[L]{GNDEC}  
\fancyfoot[R]{URN: 2435157}  
\renewcommand{\footrulewidth}{0.4pt}  
\begin{document} 
\title{Microbial Cell Surfaces}  
\author{SAKSHAM}  
\date{}  
\maketitle  
\thispagestyle{fancy}   
\section{Introduction}  
Understanding the functions of microbial cell surfaces requires knowledge of their 
structural and physical properties\cite{dufrene2002atomic}. Electron microscopy has 
long been recognized as a key technique in microbiology to elucidate cell surface ultra 
structure\cite{engel1999atomic}. An exciting achievement has been the development 
of cryotechniques, which allow high-resolution imaging of cell structures in conditions 
close to the native state\cite{franz2008atomic}. Yet direct observation in aqueous 
solution remained impossible.  
\section{Methods and Techniques}  
Because of the small size of microorganisms, the physical properties of their surfaces 
have been difficult to study\cite{marrese2017atomic}. Quantitative and qualitative 
information on physical properties can be obtained by electron microscopy techniques, 
X-ray photoelectron spectroscopy, infrared spectroscopy, contact, and electrophoretic 
mobility 
measurements 
\cite{altman2015noncontact}. 
\cite{author2023techniques,researcher2022approach,scientist202 
Recentadvances 
methodology,developer2020innovation,expert2019analysis,analyst2018study} have 
expanded our understanding of these properties.  
\bibliographystyle{plain}  
\begin{thebibliography}{10}  
\bibitem{dufrene2002atomic}  
Dufrene, Y. F. (2002). Atomic force microscopy, a powerful tool in microbiology. Journal 
of Bacteriology, 184(19), 5205-5213.  
\bibitem{engel1999atomic}  
Engel, A., & Muller, D. J. (1999). Observing single biomolecules at work with the atomic 
force microscope. Nature Structural Biology, 6(10), 811-817.  
\bibitem{franz2008atomic}  
Franz, C. M., & Muller, D. J. (2008). Analyzing focal adhesion structure by atomic force 
microscopy. Journal of Cell Science, 121(11), 1771-1780.  
\bibitem{marrese2017atomic}  
Marrese, M., Guarino, V., & Ambrosio, L. (2017). Atomic force microscopy: a powerful 
tool in biomedicine. Biomedicine & Pharmacotherapy, 93, 1023-1036.  
\bibitem{altman2015noncontact}  
Altman, M., Lee, J. R., & Tranchina, D. (2015). Noncontact atomic force microscopy: 
a tool for biological imaging. Science Advances, 1(8), e1500189.  
\bibitem{author2023techniques}  
Author, A. (2023). Advanced Techniques in Microbial Imaging. Journal of Advanced 
Techniques, 45(12), 123-134.  
\bibitem{researcher2022approach}  
Researcher, B. (2022). New Approaches to Microbial Surface Studies. Microbial 
Research Journal, 38(7), 678-689.  
\bibitem{scientist2021methodology}  
Scientist, C. (2021). Methodology in Microbial Surface Analysis. Analytical Methods, 
55(9), 345-356.  
\bibitem{developer2020innovation}  
Developer, D. (2020). Innovations in Microscopy for Microbiology. Microscopy 
Innovations Journal, 30(5), 567-578.  
\bibitem{expert2019analysis}  
Expert, E. (2019). Analysis of Microbial Surfaces: New Insights. Microbial Analysis, 
24(4), 234-245.  
\bibitem{analyst2018study}  
Analyst, F. (2018). Study of Microbial Surface Properties. Journal of Microbial Studies, 
18(3), 145-157.  
\end{thebibliography}  
\end{document} 






%Experiment 12
%Develop a LaTeX script to create a simple report and article by using suitable 
%commands and formats of Major/Minor project report. 
\documentclass[a4paper,12pt]{report}  
\usepackage[utf8]{inputenc}  
\usepackage{amsmath}  
\usepackage{amsfonts}  
\usepackage{amssymb}  
\usepackage{graphicx}  
\usepackage[left=3cm,right=2cm,top=2cm,bottom=2cm]{geometry}  
\usepackage{fancyhdr}  
\usepackage{lipsum}  
% Header and footer  
\pagestyle{fancy}  
\fancyhf{}  
\fancyhead[L]{Experiment 12}  
\fancyfoot[R]{\thepage}  
\fancyfoot[L]{GNDEC}  
\begin{document}  
\begin{titlepage} 
\centering  
{\fontsize{24pt}{28pt}\selectfont \textbf{SEMINAR AND TECHNICAL REPORT  
WRITING}}\par  
\vspace{1cm}  
{\fontsize{12pt}{14pt}\selectfont {SUBMITTED IN PARTIAL FULFILMENT OF THE  
REQUIREMENTS FOR THE AWARD OF THE DEGREE OF}}\par  
\vspace{0.5cm}  
{\fontsize{14pt}{16pt}\selectfont \textbf{BACHELOR OF TECHNOLOGY}}\par  
{\fontsize{14pt}{16pt}\selectfont 
(Computer Science and Engineering)}\par     
\vspace{2cm}  
\includegraphics[width=5cm]{GNDEC.png}\par\vspace{2cm}  
\begin{center}  
\vspace{1em}   
\begin{minipage}[t]{0.46\textwidth}  
        \raggedright  
        \textbf{SUBMITTED BY}\\[0.6em]  
        SAKSHAM\\[0.4em]  
        URN: 2435157\\[0.4em]  
        Section: CSE (E1)  
      \end{minipage}  
      \hfill  
      \begin{minipage}[t]{0.46\textwidth}  
        \raggedleft  
        \textbf{SUBMITTED TO}\\[0.6em]  
        Dr. Amandeep Kaur Sohal\\[0.4em]  
        Assistant Professor (CSE)  
      \end{minipage}  
    \end{center}  
  
\vfill  
  
    {\fontsize{14pt}{16pt}\selectfont \textbf{DEPARTMENT OF COMPUTER SCIENCE 
AND ENGINEERING}}\par  
     {\fontsize{14pt}{16pt}\selectfont \textbf{GURU NANAK DEV ENGINEERING  
COLLEGE}}\par  
    {\fontsize{14pt}{16pt}\selectfont \textbf{LUDHIANA, 141006}}\par  
    \vspace{0.5cm}  
  
\end{titlepage}  
  
\begin{abstract}  
This is a brief summary of the report. It outlines the main objectives, methods, and 
conclusions of the project.  
\lipsum[1]  
\end{abstract}  
  
\tableofcontents  
\newpage  
  
\chapter{Introduction}  
\section{Background}  
\lipsum[1]  
\section{Objectives}  
\lipsum[2]  
\newpage  
\chapter{Methodology}  
\section{Data Collection}  
\lipsum[1]  
  
\section{Analysis}  
\lipsum[1]  
  
\newpage  
  
\chapter{Results}  
\section{Findings}  
\lipsum[1]  
  
\newpage  
  
\chapter{Conclusion}  
\lipsum[1]  
  
\newpage  
  
\begin{thebibliography}{9}  
\bibitem{latexcompanion}  
Leslie Lamport,  
\textit{LaTeX: A Document Preparation System}, Addison 
Wesley, 1986.  
  
\bibitem{einstein}  
Albert Einstein,  
\textit{Relativity: The Special and the General Theory}, H. 
Holt and Company, 1920.  
\end{thebibliography}  
  
\end{document}







%IEEE
\documentclass{article}
% Use the standard article class
\usepackage[a4paper, margin=1in]{geometry}
% Set page size and margins
\usepackage{multicol}
% Allows multiple columns in the document
\usepackage{fancyhdr}
% Allows headers and footers
\begin{document}
% Begin document content
\pagestyle{fancy}
% For page headers/footers
\fancyhff}
% Clear default header and footer fields
\fancyhead[L]{2020 IEEE International Conference on Artificial Intelligence and
Computer Applications (ICAICA)}
\begin{center}
% Center the title
\textbf{\LARGE Research on Expert System in Power Network Operation Ticket}
lend{center}
\setlength{\columnsep}{8pt}
% Set spacing between columns
\begin{multicols}{3}
% Set 3-column layout
\begin{center}
% First author block centered
\textbf{Yuhang Li}||
Nari Technology (State Grid Electric Power Research Institute) Co., Ltd.Il
Nanjing, Chinall
Nari-Tech Power Distribution \& Rural Electrification Branch Companyll
Nanjing, China
lend{center}
\columnbreak
\begin{center}
% Second author
\textbf{Chenxi Huang}||
Nari Technology (State Grid Electric Power Research Institute) Co., Ltd.||
Nanjing, China||
Nari-Tech Power Distribution \& Rural Electrification Branch Company||
Nanjing, China
\end{center}
\columnbreak
\begin{center}
% Third author
\textbf{Chao Lu}||
Nari Technology (State Grid Electric Power Research Institute) Co., Ltd.Il
Nanjing, China||
Nari-Tech Power Distribution \& Rural Electrification Branch Company||
Nanjing, China
\end{center}
\end{multicols}
\begin{abstract}
\textbf{\textit{Abstract-}} This paper realizes the automatic generation system of
reversing operation ticket, aiming at the problem of insufficient flexibility and
intelligence by expert system technology. By analyzing the characteristics of the power
system, the knowledge base and the inference engine which constitute the operation
ticket expert system are studied separately to design a knowledge base with a more
concise knowledge representation. The Rete algorithm was used in the study of
reasoning mechanism, and the design of the operation ticket expert system was
completed finally.
\textit{Keywords-} artificial intelligence, operation ticket, expert system, rete algorithm
\end{abstract}
\end{document}






%intro page
\documentclass[a4paper,12pt]{article}
\usepackage{graphicx}
\usepackage{setspace}
\usepackage[margin=1in]{geometry}
\usepackage{times}

\begin{document}
\thispagestyle{empty}
\begin{center}

% ---- TITLE ----
\vspace*{1cm}
\Large \textbf{SEMINAR AND TECHNICAL\\REPORT WRITING} \\
\vspace{1cm}

\small \textbf{SUBMITTED IN PARTIAL FULFILMENT OF THE REQUIREMENTS FOR THE\\AWARD OF THE DEGREE OF} \\
\vspace{1cm}

\Large \textbf{BACHELOR OF TECHNOLOGY} \\
\textbf{(Computer Science and Engineering)} \\
\vspace{2cm}

% ---- LOGO ----
\includegraphics[width=4cm]{GNDEC.png} \\
\vspace{2cm}

% ---- SUBMITTED BY / TO ----
\begin{tabular}{c c}
\textbf{SUBMITTED BY} & \textbf{SUBMITTED TO} \\[10pt]
\textbf{SAKSHAM} & Dr. Amandeep Kaur Sohal \\
URN: 2435157 & Assistant Professor (CSE) \\
Section: CSE (E1) &  \\
\end{tabular}

\vspace{2cm}

% ---- COLLEGE DETAILS ----
\textbf{DEPARTMENT OF COMPUTER SCIENCE AND\\ ENGINEERING} \\
\textbf{GURU NANAK DEV ENGINEERING COLLEGE} \\
\textbf{LUDHIANA, 141006}

\end{center}
\end{document}
